% Reader-facing integration draft.  Its finite arithmetic propositions must
% be joined to fresh complete numerical checks and the source correspondence
% before they can be asserted as verified results in the released manuscript.
% The original complete analytic bodies are preserved in analytic_sources/.
\section{Why the intermediate indices reduce to finite local comparisons}
\label{ext:local-engine}

The argument has two distinct parts. First, contact at an extremizer
forces its marginal into a finite list of neighborhoods of two-point
laws. Second, a direct comparison of the actual iid sum with its own
binomial count excludes every neighborhood. The first part uses
extremality; the second part uses only the explicit support and moment
bounds already acquired for that same marginal. This order is essential.

The two index ranges use the same analytic argument. The finite range
$42\le n\le48$ permits a stronger Gaussian support estimate and has
92 surviving probability intervals. The uniform range $n\ge47$ uses a
different support estimate and has 83 surviving intervals. An interval
table supplies simultaneous bounds on one law. Its coordinates are not
independently selected laws, and it is not an assertion that every
point of a rectangular relaxation is realizable. The complete hypotheses
for the two applications are as follows, with
$d_n=n(C_{n-1}-C_n)_+$:
\[
\begin{array}{c|c|c|c|c}
 \text{index range}&n(C_n-\CE)&d_n&C_n&C_{n-1},C_{n-2}\\\hline
 42\le n\le48&(0,.54]&\le.19&<.4226&<.4253\\
 n\ge47&(0,.7]&\le.19&<.4247&<.4253.
\end{array}
\]
Every support and acquisition claim in this section is applied under
the corresponding full row of hypotheses.

\subsection{Contact first gives bounded support}

Let $C=C_n>\CE$ and choose the standardized attained positive ratio
maximizer from Corollary~\ref{ext:ratio-attainment}. Write
$\rho=\E|X|^3$, $\beta=\E[X|X|]$, and let $H$ be its raw sum
threshold. Set $d_n=n(C_{n-1}-C_n)_+$ and
$s_n=\sqrt{n/(n-1)}$. The contact inequality and the predecessor bound
give, at every support point $t$,
\begin{equation}\label{ext:contact-support}
 0\le \mathcal H_n(H,t)+C\left\{-|t|^3+
       \frac32\rho t^2+3\beta t+\rho e_n\right\}
       +\rho s_n d_n,
 \qquad e_n=n(s_n-1)-\frac12,
\end{equation}
where
\[
 \mathcal H_n(H,t)=n^{3/2}
 \left\{\Phi\left(\frac{H-t}{\sqrt{n-1}}\right)
              -\Phi\left(\frac H{\sqrt n}\right)\right\}
 +nt\phi\left(\frac H{\sqrt n}\right)
 +\frac H2(t^2-1)\phi\left(\frac H{\sqrt n}\right).
\]
This estimate requires only a finite third moment. In particular, it
can exclude a tiny exterior atom before any fourth or higher moment
is used.

Here is the finite calculation underlying the support check. Put
$v=n-1$, $u=|t|$, $z=H/\sqrt v$, and
\[
 A=\frac{u^3-3u}{6},\quad
 B=\frac{u^4-6u^2+3}{24\sqrt v},\quad
 D_5=\frac{u^5-15u}{120v}.
\]
Taylor's theorem in the argument and the variance of the Gaussian CDF
gives
\[
 s_n^{-3}\mathcal H_n(H,\pm u)
 =\pm A(1-z^2)\phi(z)+B\phi'''(z)
       \mp D_5\phi^{(4)}(z)+R_5,
\]
\[
 |R_5|\le \frac{51u^6}{14400v^{3/2}}
       +\frac{51(3u^2-2)}{960v^{3/2}}+\frac{u}{8v^2}.
\]
Here $u\ge3.84$, the range needed for the exterior-support test, so in
particular $3u^2-2>0$. The displayed remainder follows from the
more general expansion for $u>\sqrt3$ and is uniform over all real
thresholds. Gaussian Fourier
inversion proves the derivative bounds used here. For $|z|\le1$ use
$\phi^{(4)}(z)\ge\phi(0)(3-15z^2/2)$ and complete a square; on
$|z|\ge1$ use the uniform derivative bounds. The resulting polynomial
inequalities in $u$ are checked on adjacent closed intervals from
$3.84$ to $4.5$. The exterior rays use the cubic bound
$|\mathcal H_n(H,t)|\le u^3/15+u/5+1/2$ and its derivative signs.
Under the corresponding local hypotheses stated above, the accepted
polynomial signs exclude all support with $|t|\ge3.84$.

For a finite index band $N\le n\le L$, the favorable term
$-3D_5$ is evaluated at $L$, and the adverse terms at $N$.
For an infinite band its lower bound is zero. Thus finite support
estimates must keep their upper endpoint. The finite campaign uses
$(N,L)=(42,48)$. The uniform campaign covers
\[
 [45,49],\ [50,59],\ [60,89],\ [90,134],\ [135,199],\ [200,\infty),
\]
using only indices at least 47. These bands have no uncovered integer.
The full support polynomials and their derivation are retained in the
support supplement; no asymptotic remainder is used.

\subsection{A single nonnegative resource controls distance from two points}

Now all moments exist. For the closed support gaps $g=(a,b)$ put
$\tau_g=\E[X\mathbf1_{X\ge t}]$, constant for $a<t<b$, and define
\[
 J=\sum_g\tau_g(b-a)^2,\qquad
 D=(3+\sqrt{10})\rho-\E X^3-3J.
\]
Let $\nu$ be the law of $H-S_{n-1}$ and
\[
 L=\sum_{g=(a,b)}\tau_g\int_{[a,b)}(t-a)\,\nu(dt)\ge0.
\]
The half-open convention fixes the contribution of every atom. Contact
and compensated inversion of the smooth function $(x-z)|x-z|$ give
the exact identity
\begin{equation}\label{ext:resource-identity}
 n(C-\CE)\rho+n\CE\rho(1-e^{-w})
       +\frac{n\phi(0)e^{-w}}6D+n^{3/2}L=\mathcal R,
 \qquad w=H^2/(2n).
\end{equation}
For completeness, the right side is the following explicit expression.
Set $\gamma=\E X^3$, $M_j=\E|X|^j$ and
$\varphi=\phi(H/\sqrt n)$, and define
\[
 J_1=\sum_{g=(a,b)}\tau_g\int_a^b(t-a)t\,dt,\qquad
 J_2=\sum_{g=(a,b)}\tau_g\int_a^b(t-a)t|t|\,dt,
\]
\[
 K_1=\E[X|X|^3]-3\beta-\tfrac32\rho\gamma,\qquad
 K_q=M_5/4-\rho M_4/2-\tfrac32\beta\gamma+\rho/2,
 \qquad U=J_1+(M_4-3)/6.
\]
Then
\begin{align}\label{ext:resource-remainder}
 \mathcal R={}&C\{3\beta J/2+3\rho J_1-3J_2-K_q\}
      +H(CK_1+\varphi U)\notag\\
 &-H^2(C\rho+\gamma\varphi/2)+E_n.
\end{align}
If $f(t)=\E e^{itX}$ and
$V_n(t)=f(t)^n-e^{-nt^2/2}(1-in\gamma t^3/6)$, the exact error is
\[
 E_n=\frac{2\sqrt n}{\pi}\int_0^\infty
             \frac{\Im(e^{-itH}V_n(t))}{t^3}\,dt.
\]
The integral is absolutely convergent. A complete Fourier estimate
gives a coherent bound $|E_n|\le K_F+s_F|H|$; a constant and a slope
from unrelated estimates may not be combined. The low-frequency part retains the signed
quartic/fifth-order cancellation. The high-frequency part includes the
whole half-line and uses an event under the product of the same
marginal law. The detailed formulas appear in the Fourier supplement.
There is no inversion of a discontinuous CDF and no hidden lattice
midpoint convention.

The key geometric inequality is elementary. If the central support
gap is $(-A_0,B_0)$, let $p=\Pp(X>0)$, $q=1-p$,
$m=\E X_+=\E X_-$, $V=\E[X\mid\mathbf1_{X>0}]$ and
$\theta=\E(X-V)^2$. Then $m^2=(1-\theta)pq$. Put
$\Delta_A=m/q-A_0$ and $\Delta_B=m/p-B_0$ and
$p_*=(4-\sqrt{10})/2$. Direct conditional expansion gives
\begin{align}\label{ext:joint-gap}
 D\ge{}&\frac{(6+2\sqrt{10})(1-\theta)^{3/2}}{\sqrt{pq}}
                      (p-p_*)^2\notag\\
 &+6m\{(A_0+B_0)(\Delta_A+\Delta_B)+\Delta_A\Delta_B\}\notag\\
 &+3(\sqrt{10}+1)B_0\theta_+
       +3(\sqrt{10}+3)A_0\theta_-,
\end{align}
where $\theta_\pm=\E[(X-V)^2;\pm X>0]$.
Every term is nonnegative. If zero belongs to the support, instead
$D\ge(\sqrt{10}-1)\rho$. Consequently a sufficiently small bound on
$D$ proves that the central gap exists, then simultaneously bounds the
probability parameter, conditional noise and inward support displacement.

A first raw bound $\mathcal R<r<.4n$ licenses removal of the Gaussian
weight in \eqref{ext:resource-identity}. Indeed, before making that
change, its rearrangement proves
$\phi(H/\sqrt n)(\E X^3+3J)/6>0$. The resulting resource includes
\[
 n(C-\CE)\rho+\frac{997}{15000}nD+n^{3/2}L
 \le\min\left\{T,\max_j(A_j+F_{j,\pm}|H|-Q_jH^2)\right\}.
\]
The four corner quadratics retain the common score/moment alternatives.
The favorable Gaussian phase term is already charged in $Q_j$ and
must not be charged again.

Each acquisition step intersects previously established enclosures with
one of these inequalities. A split preserves the union of the closed
cells; an empty necessary interval excludes the cell. A support
improvement is proved before its radius enters a later moment estimate.
Induction over this finite schedule establishes all exported bounds.
No convergence of an iteration or existence of a fixed point is needed.
In particular, the weighted inequality $nD\le Nd_D$ is retained;
the unweighted inequality $D\le d_D$ alone would not bound $n\theta$.

\subsection{The comparison uses the actual binomial count}

On each surviving row write
\begin{gather*}
 X=\ell(B+U-p),\qquad \Pp(B=1)=p,\qquad \E(U\mid B)=0,\\
 v=p(1-p),\qquad s=\E U^2,\qquad\lambda=ns,
 \qquad\ell=(v+s)^{-1/2}.
\end{gather*}
The conditional noises obey
\[
 |U_j|\le A,\qquad U_0\le a_0,\qquad U_1\ge-a_1,
 \qquad d=a_0+a_1<\min(p,1-p).
\]
The row supplies these quantities and both signed threshold ranges.
Choose the nearest binomial count so that
$np+H/\ell=k+u$, $|u|\le1/2$, and put $\delta=k-np$.
The acquired ranges ensure $1\le k\le n-1$.
Let $\pi_j=\binom njp^j(1-p)^{n-j}$ and let $W_j$ be the sum of
$j$ independent copies of $U_1$ and $n-j$ independent copies of $U_0$.
Then the exact difference from the Bernoulli CDF is
\begin{equation}\label{ext:local-count-identity}
 Q=-\sum_{i=0}^k\pi_{k-i}\Pp(W_{k-i}>u+i)
       +\sum_{i=1}^{n-k}\pi_{k+i}\Pp(W_{k+i}\le u-i).
\end{equation}
Notice the strict outgoing and closed incoming events. This is simply
conditioning on the actual iid labels, so it holds at atom collisions.

Put $\rho_B=(1-2v)/\sqrt v$ and
\[
 C_B^+=\frac{\sqrt n}{\rho_B}
       \{\Pp(\operatorname{Bin}(n,p)\le k)-\Phi(\delta/\sqrt{nv})\},
\]
\[
 f(u)=\frac{\Phi((\delta+u)/\sqrt{n(v+s)})
                  -\Phi(\delta/\sqrt{n(v+s)})}{\pi_k}.
\]
Subtracting the two normalized scores yields
\begin{equation}\label{ext:comparison-master}
 \rho_B(C-C_B^+)=\sqrt n\pi_k\{Q/\pi_k-f(u)\}+\mathcal R_c,
\end{equation}
\[
 \mathcal R_c=\sqrt n\{\Phi(\delta/\sqrt{nv})
                   -\Phi(\delta/\sqrt{n(v+s)})\}+C(\rho_B-\rho).
\]
The count $k$ and $\delta$ are held fixed when differentiating $f$.
This convention matters in the signed Gaussian refinement.

\begin{lemma}[Two sufficient local tests]\label{ext:two-local-tests}
Suppose $\sqrt n\pi_k\ge m>0$ and $\lambda>0$.
The following are sufficient to prove
$\rho_B(C-C_B^+)<-\varepsilon\lambda$.
\begin{enumerate}
\item Bounds $O,T,R_0,\kappa$ satisfy
\[
 \Pp(W_k>u)+f(u)\ge O,\quad
 \sum_{i\ge1}\Pp(W_{k+i}\le u-i)\le T,\quad
 \pi_{k+i}/\pi_k\le R_0,\quad
 \mathcal R_c\le\kappa\lambda/n.
\]
If $D_c=O-R_0T>0$ and $mD_c-\kappa\lambda/n>
\varepsilon\lambda$, the conclusion holds.
\item Valid functions $B(\lambda)<0$, $I(\lambda)\ge0$ obey
\[
 \rho_B(C-C_B^+)\le\sqrt n\pi_k\lambda B(\lambda)
                         +\lambda I(\lambda)/n.
\]
For $n\ge N$, it is sufficient that
$mB(\lambda)+I(\lambda)/N<-\varepsilon$.
\end{enumerate}
\end{lemma}
\begin{proof}
For the first test discard only outgoing terms with $i>0$ in
\eqref{ext:local-count-identity}; the discarded terms are nonpositive.
Use \eqref{ext:comparison-master}, then the positive lower mass bound.
For the second test use $B<0$ before replacing $\sqrt n\pi_k$ by $m$,
and $I\ge0$ before replacing $n$ by $N$.
\end{proof}

The nearest-count estimates use a quartic minorant of the strict event,
Cantelli's inequality and a Poisson/Bennett upper-tail comparison.
For vanishing noise the quartic is rescaled by the actual variance;
its expression extends to zero noise. Thus the first interval is
$(0,.05]$, with no missing positive lower cutoff.

The whole-count estimate couples all conditional noises independently
and defines $K=\max\{j:j+W_j\le k+u\}$. Since every increment of
$j+W_j$ is at least $1-d>1/2$, this is a genuine random cutoff and
the actual sum probability is $\E F_B(K)$. Binomial curvature,
marked crossings and the nonzero fractional-part remainder give the
second test. The bound on marked crossings uses one common deleted
background law; it does not maximize two conditional probabilities
using unrelated marginals. A positive harmonic gain is obtained from
\[
 1-\cos(2\pi x)\le\frac{4553}{1000}x\quad(0\le x\le1)
\]
in the finite refinement, and its weaker constant $23/5$ in the uniform
campaign. All resulting error bounds are proportional to $\lambda$.
This is why the comparison remains valid for arbitrarily small defects.

The joint inequality \eqref{ext:joint-gap} is also retained in the local
comparison. In the finite refinement it gives a resource of the form
$d/d_*+\lambda/\Lambda_*\le1$. On a noise interval
$\lambda\ge L_0$, one may therefore use
$d\le d_*(1-L_0/\Lambda_*)$. The two separate ceilings would lose
the information needed near the lower end of the probability range.

\subsection{The two complete acquisitions and their phase covers}

The finite campaign starts with 100 closed probability cells. Refining
the surviving cells gives the following partition, including all seams:
\[
\begin{array}{c|c}
 {[0,.23],\ [.63,1]}&\text{initial exclusion}\\
 {[.23,.27]}&4\text{ coarse exclusions}\\
 {[.27,.29]}&10\text{ fine exclusions}\\
 {[.29,.38]}&45\text{ left rows}\\
 {[.38,.474]}&47\text{ central rows}\\
 {[.474,.48]}&3\text{ central exclusions}\\
 {[.48,.63]}&15\text{ upper exclusions}.
\end{array}
\]
There are 184 closed cells altogether. Endpoint probabilities zero and
one cannot occur for a centered variance-one law. The left support
radius is first contracted to $3.025$ and then row by row. Contact
proves the strip $(-.575,1.575)$ in Bernoulli coordinates before
concentration supplies the conditional-noise amplitude. The ensuing
feedback and four noise-dependent profiles use already acquired data.

On the left, the local witnesses have the exact domain cover
\[
\begin{array}{c|c|c}
 \lambda&\text{nearest-count test}&\text{whole-count test}\\
 (0,.05]&\text{all }u&--\\
 {[.05,.08]}&\text{all }u&--\\
 {[.08,.10]}&|u|\ge.19&|u|\le.19\\
 {[.10,.12]}&|u|\ge.27&|u|\le.27\\
 {[.12,\Lambda]}&--&\text{all }u.
\end{array}
\]
The central rows have complementary phase radii $.09,.14,.19,.29$
on the successive bands $[.05,.06]$, $[.06,.08]$, $[.08,.10]$,
$[.10,.12]$; beyond $.12$ the whole-count test covers all phases.
Every interval is intersected with the row's actual $\Lambda$, and
both signs of $H$ are present. The amplitude profiles on
$[.05,.08]$, $[.08,.10]$, $[.10,.12]$, $[.12,\Lambda]$ are
$.40,.35,.30,.275$ on the left and $.375,.325,.275,.25$ centrally,
each intersected with the previously proved global amplitude.

The uniform campaign starts with radius $3.84$ and 13 signed resource
passes. Central and complementary acquisition are followed by support
contraction and concentration; the central antichain step consumes the
original central data, then intersects the first improved table. This
prevents a circular use of its own output. A signed resource step
excludes one left probability interval. The final cover has 38 left
rows on $[.304,.38]$ and 45 central rows on $[.38,.47]$.
All complementary probabilities have already been excluded.
The initialization stores a floor $N=47$ and an old upper-index field
$98$. The latter is used only to check that the stored band is ordered
and in the lower-score expression $.4095+K/98$, with $K=0$ here.
It therefore imposes no upper bound on the actual index. The analytic
licenses still matter: the all-threshold support cover ends in an
infinite band, the bounded-threshold series controls the entire interval
$0\le1/n\le1/47$, the Fourier bounds have explicit monotonicity gates,
and the local tests retain the signs needed before replacing $n$ by 47.
Together these facts, rather than the name of a table, establish the
uniform index scope.

Its nearest-count comparison covers $(0,.05]$ completely. The central
MID phase radii on the next four bands are $.09,.13,.18,.29$; the left
radii are $.09,.14,.19,.29$. On the final left band the positive
phases $[.29,.5]$ use a further whole-count test. The central,
phase-resolved left and extra positive-left whole-count objects cover
the complementary phases and every phase above $.12$. All seams are
closed. These uniform constants are not substituted into the finite
campaign, or conversely.

The specified target for both campaigns is the strict local inequality
with $\varepsilon=.001$ on every remaining positive-noise cell.
Each calculation is a finite
collection of rational inequalities, with radicals and transcendental
terms replaced by proved outward rational enclosures. The checkers
test all required signs before division or multiplication, all rectangle
corners, and complete probability/noise/phase coverage.

At $\lambda=0$, conditional centering implies $U=0$ almost surely.
The marginal is Bernoulli and the all-probability finite-$n$ Bernoulli
theorem gives $C\le\CE$. For $\lambda>0$ the same theorem gives
$C_B^+\le\CE$, while Lemma~\ref{ext:two-local-tests} gives
$C<C_B^+$. Thus each acquired neighborhood contradicts $C>\CE$.
This completes the mathematical reduction of the two local exclusions
to their specified finite arithmetic predicates and the Bernoulli
theorem. The executable correspondence must retain those precise
predicates and scopes.
